% !PRINTED 481
\weberpage{481}

And likewise:

6. \gesperrt{If in an invariant of $G$ one replaces the powers and products of the variables by the corresponding derivatives of a contravariant, the result is again a contravariant.}

\begin{center}
\S. 130.

\gesperrt{Solution of the equation of the seventh degree with the group $P_{168}$ by the form problem of the group $G_{168}$.}
\end{center}

The linear substitution group $G_{168}$ has the remarkable property of being contragredient to itself. For the generating substitutions $\tau$, $\omega$ of $G_{168}$ (\S.\,115) remain unchanged under transposition, and hence if one transposes any substitution of the group, one obtains a substitution which likewise occurs in the group.

The contravariants of $G_{168}$ are therefore (apart from the designation of the variables) identical with its invariants.

But the octahedral group $G_{24}$ contained in the group $G_{168}$, which consists of the substitutions $\chi^\lambda\omega^\mu\Theta^\nu$, is distinct from its contragredient group; for by \S.\,72, (17) we have, e.\,g.,
\[
\Theta^3 = \omega\tau\omega\tau^6,
\]
and the substitution transposed to it,
\[
\Theta_1^3 = \tau^6\omega\tau\omega
\]
can, by \S.\,72, (15), be brought into the form $\tau\chi^2\Theta^2$ and is therefore not contained in $G_{24}$.

Now let $\eta_0$ be any function of the quantities $\lambda_0, \lambda_1, \ldots, \lambda_6$ belonging to the group $P_{24}$ (\S.\,128), e.\,g. the root $\lambda_0$ itself. Under the cyclic permutations $\tau^\nu$ let $\eta_0$ pass into $\eta_0, \eta_1, \ldots, \eta_6$. In the resolvents
\weq{1}{\tilde\omega_r = \sum_{0,6}^{\nu} \varepsilon^{\nu r}\eta_\nu \qquad r = 1, 2, \ldots, 6}
we then have a system of functions which under the permutations of $P_{168}$ substitute linearly, though not ternarily; since one can indeed express the $\eta_\nu$ themselves linearly in terms of the $\tilde\omega_r$.
